\documentclass[11pt]{article}
\usepackage[a4paper,margin=25mm]{geometry}
\usepackage[T1]{fontenc}
\usepackage{lmodern,amsmath,amssymb,amsthm,mathtools,longtable,hyperref}
\setlength{\emergencystretch}{2em}
\newtheorem{proposition}{Proposition}[section]
\newtheorem{corollary}[proposition]{Corollary}
\newcommand{\Ric}{\operatorname{Ric}}
\newcommand{\tr}{\operatorname{tr}}
\newcommand{\Id}{\operatorname{Id}}
\newcommand{\RR}{\mathcal R_+}
\title{Negative Goldberg: Joining the Torsion-Zero and Green-Potential Arguments}
\author{Research continuation and cross-conversation synthesis}
\date{15 October 2026}
\begin{document}
\maketitle
\begin{abstract}
We join five retrieved research notes from the conversation entitled
``Yau 26 23 19'' with the two torsion-zero continuations of 15 October.
The joined identities characterize the intersection of the torsion-zero
set with the double-maximal-eigenvalue locus: it consists exactly of the
zeros at which the second covariant derivative of the fundamental form
also vanishes. Every closed smooth torsion-zero curve with dense rank-two
locus must meet this intersection. A Green identity retains the global
curvature term at these points; its vanishing does not follow from their
finite-order local vanishing. The unrestricted negative Goldberg
conjecture remains open in these arguments.
\end{abstract}

\section{The common problem and conventions}
Let $(M^4,g,J,\Omega)$ be smooth, closed, connected and almost-K\"ahler,
oriented by $J$, with
\[
\Ric=-ag,\qquad a>0,\qquad s=-4a,\qquad |\Omega|^2=2.
\]
The target is $\nabla\Omega=0$, without additional curvature-sign or
zero-set assumptions, or an actual closed counterexample.
Let $\Delta=-\operatorname{div}\nabla$, $n=\Omega/\sqrt2$, and
\[
q=|\nabla\Omega|^2,\quad Z=\{q=0\},\quad
\RR=W^+-\frac a3\Id,\quad A=-\RR,\quad
U=\sigma_2(A)=\frac{a^2}{3}-\frac12|W^+|^2.
\]
Operator norms are Hilbert--Schmidt norms. Relative to
$\mathbb R n\oplus n^\perp$, write
\[
W^+=\begin{pmatrix}t&v^*\\v&-t\Id/2+B\end{pmatrix},\qquad
t=-\frac{2a}{3}+\frac q4,\qquad v=\eta/\sqrt2,\qquad\tr B=0.
\]
Here $\eta=W^+(\Omega)-t\Omega$ and $B$ is the trace-free two-dimensional
block, customarily denoted $W_3^+$. The scalar $a$ is not the torsion
one-form used in earlier notes. We denote that one-form by $\vartheta$:
\[
\nabla\Omega=\vartheta\otimes\phi-J\vartheta\otimes J\phi,
\qquad q=4|\vartheta|^2,
\]
where $|\phi|^2=|J\phi|^2=2$.

\section{Source correspondence and retained results}
The following five complete source texts were retrieved and read.
They are records of the other conversation, not independent publications.
\begin{longtable}{p{.44\textwidth}p{.50\textwidth}}
\textbf{Source and date}&\textbf{Mathematical content retained}\\\hline
\path{Q4_Torsion_Zero_Structure.tex}, 9 October
&Second-derivative identity; complete first-derivative normal form;
positive transverse indices; the scalar current retained before integration.\\[4pt]
\path{Q4_Canonical_Orthogonality.tex}, 12 October
&Orthogonality of $\eta$ to closed real $J$-anti-invariant forms;
positive integrating-factor criterion and its noncompact local obstruction.\\[4pt]
\path{Q4_Harmonic_Product_Reduction.tex}, 13 October
&Fixed-metric harmonic multiplication kernel, its exact rank-one constraint,
and fourth-order norm-variance behavior near a K\"ahler direction.\\[4pt]
\path{Q4_Determinant_Sign_Rigidity.tex}, 13 October
&Rigidity under $\det\RR\le0$ everywhere, from the Einstein Weyl equation
and existence of a torsion zero.\\[4pt]
\path{Q4_Green_Potential_Rigidity.tex}, 14 October
&Rigidity if the Green potential of $\det\RR$ is nonpositive at one
torsion zero; exact weighted determinant balance.\\\hline
\end{longtable}
The current conversation contributes the two 15 October notes
\path{Goldberg_Torsion_Zeros_and_Weyl_Equation_20261015.tex}
and \path{Goldberg_Global_Connections_20261015.tex}.
Their relevant statements are the exclusion of compact clean zero curves,
the normal derivative equation on a smooth zero curve, and the full
coupled Weyl equation. The zero-index calculation overlaps the 9 October
source and is counted once.

The source titled \path{Q4_Complete_Example_and_Fino_Question.tex}
was located by title but not re-read for this synthesis. No additional
claim about that construction or historical priority is made here.
This record covers the retrieved mathematical notes, not every message
of the other conversation.

For orientation, the canonical orthogonality result reads
\[
\int_M\langle\eta,\alpha\rangle\,dV=0
\quad(d\alpha=0,\ \alpha\in\Gamma(\Lambda_J^-)).
\]
The source uses the canonical-divisor theorem and the canonical curvature
formula. Thus a globally smooth $f>0$ with $d(f\eta)=0$ implies $\eta=0$.
The compact third-Gray theorem of Apostolov--Armstrong--Dr\u aghici
\cite{AAD} then implies K\"ahlerness. Einstein Ricci is $J$-invariant,
and $\eta=0$ says precisely that $\Omega$ is a Weyl eigenform;
these are the two conditions in that theorem. Existence of the multiplier
is not supplied by this implication.

\section{The common global identity, including the local zero jet}
The Einstein Weyl equation is
\[
\nabla^*\nabla W^++\frac s2 W^+-6(W^+)^2+2|W^+|^2\Id=0.
\]
Pairing it with $W^+$ and shifting by $a\Id/3$ gives
\begin{equation}\label{scalar}
(\Delta+2a)U=|\nabla A|^2-18\det\RR.
\end{equation}
Indeed,
\[
\Delta|W^+|^2=-s|W^+|^2+36\det W^+-2|\nabla W^+|^2,
\]
\[
|W^+|^2=\frac{2a^2}{3}-2U,
\qquad \det W^+=-\frac{2a^3}{27}+\frac a3U-\det A.
\]
Define
\[
\Psi=(\Delta+2a)^{-1}\det\RR,\qquad
\mathcal D=(\Delta+2a)^{-1}|\nabla A|^2.
\]
The inverse exists on all scalar functions, including constants, and
has a strictly positive Green kernel $G_a$ on a connected closed manifold.
Thus
\begin{equation}\label{resolvent}
U=\mathcal D-18\Psi,\qquad \mathcal D\ge0.
\end{equation}

At $p\in Z$, the harmonic-form equation gives
$\langle An,n\rangle=a$, so
\[
A(p)=\begin{pmatrix}a&z^*\\z&C\end{pmatrix},\qquad
\tr C=0,\quad z=-v,\quad C=-B.
\]
Consequently
\begin{equation}\label{zerou}
U(p)=-|z|^2-\tfrac12|C|^2
     =-\tfrac12(|\eta|^2+|B|^2)(p).
\end{equation}
The local torsion-zero identity, established in the 9 October note and
retained in the present conversation, is
\begin{equation}\label{hessian}
|\nabla^2\Omega|^2=4|B|^2-6|\eta|^2
\qquad\text{on }Z.
\end{equation}
It follows by twice differentiating
$t=\langle W^+n,n\rangle=s/6+q/4$ at $\nabla n=0$:
\[
\Delta t=3|\eta|^2-2|B|^2,
\qquad \Delta q=-2|\nabla^2\Omega|^2.
\]

\begin{proposition}[Joined zero-point identity]\label{joined}
Every $p\in Z$ satisfies
\begin{equation}\label{master}
\boxed{18\Psi(p)=
\int_MG_a(p,y)|\nabla A(y)|^2\,dV_y
+\frac18|\nabla^2\Omega(p)|^2+\frac54|\eta(p)|^2.}
\end{equation}
\end{proposition}
\begin{proof}
Combine \eqref{resolvent}--\eqref{hessian}; specifically
$|B|^2=|\nabla^2\Omega|^2/4+3|\eta|^2/2$.
\end{proof}
This is a direct connection between the two existing arguments, not a
second independent global rigidity theorem. The Green theorem is recovered:
if $\Psi(p)\le0$ at one torsion zero, all terms on the right vanish.
Positivity of $G_a$ forces $\nabla A=0$ globally. Then $U=\det A=0$,
$\tr A=a$, and $A$ is a parallel rank-one operator with eigenvalues
$(a,0,0)$. Set $\Pi=A/a$. The harmonic-form equation gives
\[
\nabla^*\nabla(\Pi n)=2\Pi(a\Id-A)n=0.
\]
Hence $\Pi n$ is parallel, and at $p$ it has length one. Since $|n|=1$,
this gives $n=\Pi n$ globally, hence $\nabla\Omega=0$.

\section{What the spectral intersection means exactly}
Order the Weyl eigenvalues by $\lambda_1\ge\lambda_2\ge\lambda_3$ and set
\[
\Sigma=\{\lambda_1(W^+)=\lambda_2(W^+)\}.
\]
This definition uses the \emph{largest} repeated eigenvalue.

\begin{proposition}\label{intersection}
On the torsion-zero set,
\[
\boxed{Z\cap\Sigma
=\{p:\nabla\Omega(p)=0,\ \nabla^2\Omega(p)=0\}.}
\]
At every such point, $\eta=B=0$ and $\RR=-a\,n\otimes n$.
\end{proposition}
\begin{proof}
Suppose $p\in Z\cap\Sigma$. Let $\lambda$ be the repeated maximal
eigenvalue of $\RR$; its remaining eigenvalue is $-a-2\lambda$.
For a unit vector $e$ in the remaining eigenline,
\[
\RR=\lambda\Id-(a+3\lambda)e\otimes e.
\]
The Rayleigh value at a torsion zero is $\langle\RR n,n\rangle=-a$.
The operator cannot be scalar there. Thus
\[
|\langle e,n\rangle|^2=\frac{a+\lambda}{a+3\lambda},\qquad
\lambda\ge0.
\]
The last inequality also follows from the smallest-eigenvalue bound
$-a\ge-a-2\lambda$. Computing the off-diagonal block and the traceless
restriction to $n^\perp$ gives
\[
|\eta|^2=4\lambda(a+\lambda),\qquad |B|^2=2\lambda^2.
\]
Equation \eqref{hessian} becomes
\[
|\nabla^2\Omega|^2=-24a\lambda-16\lambda^2.
\]
Nonnegativity forces $\lambda=0$ and $\nabla^2\Omega=0$, and the
Rayleigh equality then gives $e=\pm n$.

Conversely, at $p\in Z$ with $\nabla^2\Omega=0$, the harmonic-form
equation $\nabla^*\nabla\Omega=(q/2)\Omega+2\eta$ gives $\eta=0$.
Equation \eqref{hessian} gives $B=0$. The eigenvalues of $W^+$ are
$(-2a/3,a/3,a/3)$, proving $p\in\Sigma$.
\end{proof}

At the intersection, \eqref{master} reads
\begin{equation}\label{intersectiongreen}
\boxed{18\Psi(p)=\int_MG_a(p,y)|\nabla A(y)|^2\,dV_y.}
\end{equation}
It does not imply $\Psi(p)\le0$. If $\nabla A$ is nonzero somewhere,
the right side is strictly positive. Thus ``the zero set meets the
spectral locus'' and ``the Green rigidity hypothesis holds'' are separate
statements. Vanishing of two derivatives is also not infinite-order
vanishing for unique continuation.

\section{A geometric connection supplied by closed zero curves}
Let $K=\Lambda^{2,0}T^*M$ and regard torsion as a section
$\sigma\in\Gamma(T^{1,0}M\otimes\bar K)$, with the normalizations of
the zero-structure notes. The canonical Hermitian connection satisfies
\[
D_X\sigma=X\otimes c,\qquad X\in T^{1,0}M.
\]
Let $C$ be a closed smooth $J$-holomorphic curve contained in $Z$, with
dense rank-two locus for $D\sigma$. The current conversation establishes
that $C$ is totally geodesic and that its normal torsion derivative is a
smooth section
\[
\zeta\in\Gamma\bigl(\operatorname{Hom}_{\mathbb C}
 (\bar\nu_C,T^{1,0}C\otimes\bar K)\bigr)
\simeq\Gamma(\bar K^{\otimes2}|_C).
\]
Here $\zeta$ is unrelated to the off-diagonal Weyl notation in some
earlier manuscripts. On $C$, all tangent derivatives and all $(1,0)$
derivatives of $\sigma$ vanish; $\zeta$ is the remaining normal
$(0,1)$ derivative.

\begin{corollary}\label{curve}
Every such $C$ meets $Z\cap\Sigma$.
\end{corollary}
\begin{proof}
The canonical curvature formula restricted to a torsion-zero curve gives
\[
c_1(TM,J)[C]=\frac{s}{8\pi}\operatorname{Area}(C).
\]
As $c_1(\bar K)=c_1(TM,J)$,
\[
\deg(\bar K^{\otimes2}|_C)=\frac{s}{4\pi}\operatorname{Area}(C)<0.
\]
A nowhere-zero smooth section would trivialize this complex line bundle
and force degree zero. Hence $\zeta(p)=0$ somewhere. All of $D\sigma$
then vanishes at $p$. Since $\sigma(p)=0$, the canonical and
Levi-Civita derivatives agree there, and the torsion identification gives
$\nabla^2\Omega(p)=0$. Proposition \ref{intersection} applies.
\end{proof}
This corollary requires neither isolated zeros of $\zeta$ nor a
positivity-of-indices assertion. It extends the earlier clean-curve
exclusion by identifying where a rank drop must occur spectrally.
It does not assert that every component of the torsion-zero set is a
smooth curve.

The equation retained on $C$ is
$D^{1,0}\zeta=r$, with $r$ a normal derivative of $c$, and
$D^{1,0}r=0$ under the stated dense rank-two hypothesis.
It does not say that $\zeta$ is a closed $J$-anti-invariant two-form
on $M$. The canonical-divisor theorem used in the orthogonality note
therefore cannot be applied to this section by changing its name.

\section{The resulting common continuation point}
The merged statements are compatible and give the following precise
restrictions on any non-K\"ahler closed candidate:
\begin{enumerate}
\item $Z\ne\varnothing$, by the classical Nijenhuis-bundle obstruction
$c_2(T^{1,0}M\otimes\bar K)[M]=5\chi+6\tau>0$.
\item $\Psi(p)>0$ at every $p\in Z$, with its exact value prescribed by
\eqref{master}. In particular $\det\RR>0$ somewhere.
\item At a torsion zero, membership in $\Sigma$ is exactly second-order
torsion vanishing. Any closed smooth zero curve of the preceding type
must contain such a point.
\item Neither that intersection nor the finite-order vanishing removes
the global integral in \eqref{intersectiongreen}.
\end{enumerate}
Accordingly, the direct unresolved connection on this route is a
restriction on the actual Einstein--almost-K\"ahler system that forces
the nonnegative integral in \eqref{intersectiongreen} to vanish, or that
forces the Green sign at a torsion zero without a zero-curve hypothesis.
The present identities do not furnish that restriction. A valid proof
may use a different route; none of these auxiliary conditions is added
to the original conjecture.

The harmonic multiplication formulation remains an exact alternative
for a fixed metric, not an additional curvature inequality. Likewise,
the existing full $B$ equation retains first-order coupling to $\eta$;
it cannot be replaced by a homogeneous positive-potential equation.
These observations preserve both conversations' usable results while
keeping the remaining implication mathematically explicit.

\begin{thebibliography}{9}
\bibitem{AAD} V.~Apostolov, J.~Armstrong and T.~Dr\u aghici,
\emph{Local models and integrability of certain almost K\"ahler
4-manifolds}, Math. Ann. 323 (2002), 633--666.
\url{https://arxiv.org/abs/math/0109098}. Lemma 1 and Theorem 2.
\bibitem{Derdzinski} A.~Derdzi\'nski,
\emph{Self-dual K\"ahler manifolds and Einstein manifolds of dimension four},
Compositio Math. 49 (1983), 405--433.
\url{https://archive.numdam.org/item/CM_1983__49_3_405_0/}.
\bibitem{LeBrun} C.~LeBrun,
\emph{Einstein metrics, symplectic minimality, and pseudo-holomorphic curves},
Ann. Global Anal. Geom. 28 (2005), 157--177.
\url{https://arxiv.org/abs/math/0409290}. Proposition 2.1.
\bibitem{Project} The seven research notes identified in Section 2.
The Green argument is inherited from the other conversation; the
closed-zero-curve argument is inherited from the present conversation.
Propositions \ref{joined} and \ref{intersection} and
Corollary \ref{curve} state their explicitly derived connections.
\end{thebibliography}
\end{document}
